\section{Globalización, Singapur y el malentendido de la «huida»}

Al final de la entrevista, Peak habla de Singapur. Su respuesta es pragmática: Manus eligió Singapur no para «huir», sino porque la empresa atiende al mercado global desde el principio y necesita una sede que le permita atender a clientes internacionales, cumplir la normativa, reunir al equipo en una misma oficina y operar a escala global. Esta pregunta tiene que ver con los productos de Agent: si los usuarios, los clientes y los pagos están repartidos por el mundo, el equipo tiene que ocuparse con anticipación del cumplimiento normativo, los datos, los pagos, los aspectos legales y la confianza de los clientes.

\subsection{Por qué no es una simple elección geográfica}

Ya se dijo que Singapur no fue una decisión emocional; esta sección desglosa las razones reales. Peak menciona que Manus ya completó trabajos de cumplimiento como SOC 2, ISO 27001 / 27701 y GDPR. Para una herramienta común dirigida al consumidor, este trabajo puede parecer excesivo, pero para una empresa de Agent es fundamental, porque un Agent puede tener acceso a cuentas de usuario, archivos, navegadores, información empresarial, pagos y operaciones automatizadas. Que los clientes de otros países estén dispuestos a usarlo suele depender de que primero confíen en su seguridad y su cumplimiento normativo.

\begin{knowledgebox}{La presión de cumplimiento normativo en una empresa global de Agent}
\begin{tabularx}{\textwidth}{p{0.22\textwidth}p{0.34\textwidth}X}
\toprule
Aspecto de cumplimiento o gobernanza & Pregunta que debe responder & Relación con el Agent \\
\midrule
Protección de datos & Dónde están los datos de los usuarios y cómo se tratan & El Agent necesita mucho contexto y muchos archivos. \\
Auditoría de seguridad & Si el sistema tiene controles y trazabilidad & Debe poder asignarse responsabilidad por cada operación automatizada. \\
Servicio transfronterizo & En qué mercado está el cliente y cómo se firma el contrato & Un producto global debe adaptarse a reglas distintas. \\
Colaboración del equipo & Si el equipo está concentrado y cuánto cuesta comunicarse & Un equipo en etapa temprana necesita iterar y decidir con rapidez. \\
\bottomrule
\end{tabularx}
\end{knowledgebox}

\subsection{Por qué «huida» es una descripción equivocada}

Hablar de «huida» confunde la globalización, el cumplimiento normativo y el mercado de clientes con una fuga geográfica. Según Peak, Manus siempre ha trabajado para el mercado global, y Singapur responde a la necesidad de atender a los clientes y organizar al equipo. En una empresa de IA, la elección de oficinas y de sede cambia según dónde estén los clientes, dónde estén los requisitos de cumplimiento y dónde estén las estructuras de pagos y legales.

\begin{warningbox}{No reducir la globalización a una palabra cargada de emoción}
Que una empresa de IA elija una sede en el extranjero puede deberse a varias razones a la vez: clientes, cumplimiento normativo, financiamiento, pagos, talento y colaboración del equipo. Resumirlo como «huida» oculta los verdaderos problemas comerciales e institucionales.
\end{warningbox}

\subsection{Resumen de la sección}

La elección de Singapur muestra que una empresa de Agent no es solo un equipo de producto en línea. Tiene que enfrentarse a clientes globales, marcos de cumplimiento normativo, seguridad de datos y colaboración organizacional. Estos problemas también forman parte del lado práctico de que «la IA se parece a la manufactura».

